% !PRINTED 377
\weberpage{377}
\begin{center}
§. 122.

Quadratic Irrational Numbers.
\end{center}

Simple and beautiful laws emerge when one applies the continued fraction expansion to the determination of the roots of a quadratic equation. The root of a quadratic equation with integral coefficients has the form
\weq{1}{\omega = x + y\sqrt{d},}
in which $x$, $y$, $d$ are integral or fractional rational numbers. We may, however, without restricting the generality, \gesperrt{assume $d$ to be an integer and without square divisor}; for if $d$ has a denominator, we can multiply through by this denominator, and may reckon the root of the square denominator, and of any square divisor of the numerator, to $y$; we further assume that $d$ is not equal to 1, since otherwise $\omega$ would be rational; $d$ may be positive or negative, and on this it depends whether $\omega$ is real or imaginary.

If we change the sign of the root, we obtain
\weq{2}{\omega' = x - y\sqrt{d},}
which is called the number \gesperrt{conjugate} to $\omega$. The product of the two numbers $\omega$, $\omega'$
\[
\omega\omega' = x^2 - y^2 d
\]
is called the \gesperrt{norm of} $\omega$ (or also of $\omega'$). \gesperrt{The norm of a product of two quadratic irrational numbers is equal to the product of the norms.}

The numbers $\omega$ and $\omega'$ are the roots of an equation with rational coefficients:
\weq{3}{\omega^2 - 2x\omega + (x^2 - y^2 d) = 0.}
To transform it into an equation with integral coefficients, we set
\weq{4}{x^2 - y^2 d = -\frac{a}{c}, \quad 2x = \frac{b}{c},}
in which $a$, $b$, $c$ are integers without common divisor, and obtain from (3)
\weq{5}{c\omega^2 = a + b\omega,}
that is, a \gesperrt{primitive} equation with integral coefficients\par % [sic] source omits terminal punctuation; the next leaf begins a new sentence
